% titel: Blatt 10, A5: Autoprotolyse des Wassers (c(H+), pH, pOH)
% skript: Mortimer Kap. 17 (Autoprotolyse, Ionenprodukt des Wassers)

% Teil I: Werkzeuge, kein Lösungsweg.

Wasser autoprotolysiert: H$_2$O$+$H$_2$O$\rightleftharpoons$H$_3$O$^+$+$OH^-$.
Das Ionenprodukt $K_w$ verknüpft in wässriger Lösung alle vier Größen.

\begin{werkzeugkasten}[pH/pOH und Ionenprodukt]
Bei 25\,°C gilt
\[
  K_w = c(\mathrm{H^+})\,c(\mathrm{OH^-}) = 10^{-14}\ mol^2L^{-2}.
\]
Und logarithmisch
\[
  pH = -\lg c(\mathrm{H^+}),\quad pOH = -\lg c(\mathrm{OH^-}),\quad
  pH + pOH = 14 .
\]
Neutral: $c(\mathrm{H^+}) = c(\mathrm{OH^-}) = \sqrt{K_w}$.
\textit{Skript:} Mortimer Kap. 17 (Ionenprodukt, pH).
\end{werkzeugkasten}

\begin{achtung}[Darauf achten]
\begin{itemize}
  \item Saure: $c(\mathrm{H^+})>10^{-7}$, basisch: $c(\mathrm{H^+})<10^{-7}$.
  \item $K_w$ ist temperaturabhängig (bei 0\,°C nicht $10^{-14}$!) --- die
        neutrale Lösung hat dann andere $c$ $-Werte$.
  \item pH-Änderung um $\Delta pH$ ändert $c(\mathrm{H^+})$ um den Faktor
        $10^{-\Delta pH}$.
  \item Ohne Fremdion-Zusatz gilt bei neutral $c(\mathrm{H^+})=c(\mathrm{OH^-})$.
\end{itemize}
\end{achtung}
